\section{The exact scope of the HNN construction}\label{sec:app-hnn}
The compatible-family HNN theorem in the common paper supplies full coset Helly and all intersections from an indexed partial bijection. Firmness turns the source's unlabelled parabolic-lattice admissibility into precisely that indexed hypothesis. It therefore recovers the whole firm-input branch of Theorem~3.35, with or without residual connectedness. Independently, deleting universal types in a residually connected input reduces its literal source construction to the firm active geometry and then adds back singleton types. Thus the flag-transitivity and residual-connectivity conclusions of part~(b) hold even for nonfirm inputs. Stable panels are infinite; the other exact panel indices are the common theorem's indices.

These two positive branches do not imply the unrestricted first sentence or part~(a) of Theorem~3.35. We give both the simpler failure of the source intersection formula and a complete counterexample to the broader conclusion. Neither is used as a substitute for the positive arguments.

\subsection*{Why an unlabelled lattice is not sufficient in general}
Let $A$ be the additive group of $\mathbf F_2^3$, with basis
$e_1,e_2,e_3$, and take four indexed subgroups
\[
 A_1=\langle e_2,e_3\rangle,\quad
 A_2=\langle e_1,e_3\rangle,\quad
 A_3=\langle e_1,e_2\rangle,\quad
 A_4=\langle e_3\rangle=A_1\cap A_2.
\]
This coset geometry is flag-transitive. Indeed, each coset prescribes
values of some coordinates. Pairwise intersecting cosets have no
conflicting coordinate prescriptions, so all of them have a common
representative. Every such flag extends to a chamber by that representative.
The geometry is not firm: $H=0$ and, for example, $R_1=0$.

Put $J=K=\varnothing$, so both associated subgroups are $A$, and let
$\phi$ cyclically permute the basis,
$e_1\mapsto e_2\mapsto e_3\mapsto e_1$.
The unlabelled standard parabolics are exactly the eight coordinate
subspaces. Thus $\phi$ preserves their whole lattice and is admissible
under the literal Definition 3.30. Its action on the three planes is
$A_1\mapsto A_2\mapsto A_3\mapsto A_1$. The image
$\phi(A_4)=\langle e_1\rangle$ is a standard parabolic but is not one
of the four individually indexed subgroups. Neither is
$\phi^{-1}(A_4)=\langle e_2\rangle$.
The equivalence classes in Definition 3.31 are consequently
$\{1,2,3\}$ and $\{4\}$.

For $s=\{4\}$ the compatibility asserted in the proof of Lemma 3.32
fails:
\[
 \phi(A_s\cap A)=\langle e_1\rangle
 \ne\langle e_3\rangle=A_s\cap A.
\]
The resulting intersection formula on source page 29 is false as well.
In the HNN extension
\[
 G=\langle A,t\mid t^{-1}vt=\phi(v)\ (v\in A)\rangle
\]
one has $G_s=\langle t,e_3\rangle=G$, since conjugating $e_3$ by
$t$ and $t^2$ supplies $e_1$ and $e_2$. Also $F=J\setminus K$ is
empty, so $G_t=A_F=A$. It follows that
\[
 G_t\cap G_s=A,\qquad A_{F\cup s}=A_4=\langle e_3\rangle,
\]
whose orders are eight and two. Normal forms ensure that $A$ embeds
in this HNN extension, so these orders do not collapse.

This example refutes the claimed compatibility and the displayed
intersection formula in their stated generality. It does not refute
the existence of the abstract HNN extension, nor does it show that
$G_s$ cannot be a special subgroup with some larger base subgroup.
It also does not disprove the source's broad flag-transitivity conclusion:
the particular output just constructed is still flag-transitive. Its
maximal parabolics are $\langle t\rangle,G,A$; the $G$-coset is a
single element incident with everything, and $A\langle t\rangle=G$
shows that every pair of the other types meets and has a common
representative.
The firm case is covered by the common firm-admissibility corollary; more general
inputs require explicit indexed compatibility or a separately proved
replacement construction.


\subsection{The exact source boundary}
Throughout, incidence means nonempty intersection of left cosets. The
source's Section 2.1, pages 6--8, permits arbitrary indexed subgroups as
the standard maximal parabolics; this name does not require maximality in
the subgroup lattice. It sets $A_J=\bigcap_{j\in J}A_j$ and
$A_\varnothing=A$. A flag-transitive coset incidence system is a geometry,
but firmness and residual connectedness are separate conditions.

Definition 3.30, pages 27--28, allows $C=A_J$ and $D=A_K$ for arbitrary
$J,K\subseteq I$. An admissible isomorphism $\phi:C\to D$ must preserve
the unlabelled lattices of standard parabolics contained in these groups,
including their intersections. In particular, $J=K=\varnothing$ and
$C=D=A$ are allowed. Theorem 3.35, page 30, assumes only flag transitivity
of the input and this admissibility for its initial assertion. Firmness
and residual connectedness occur as additional hypotheses for subsequent
conclusions. No such hidden assumptions are imposed in the definitions.

For $J=K=\varnothing$, Definition 3.31 identifies named indices whenever
$\phi(A_i)=A_j$, taking the equivalence relation generated by these pairs.
For a resulting class $s$ it prescribes
\[
 G_s=\gen{t,A_s},\qquad A_s=\bigcap_{i\in s}A_i,
 \qquad G_t=A,
 \qquad G=\gen{A,t\mid t^{-1}xt=\phi(x)\ (x\in A)}.
\]
Preservation of the unlabelled lattice does not assert that $\phi$ permutes
the individually named $A_i$. This is the distinction used here.

\subsection{An elementary input geometry}
In $A=S_{\{1,2,3,4\}}$, put
\[
 a=(34),\quad b=(12),\quad r=(13)(24),\quad s=(14)(23),
 \qquad k=ab=rs=(12)(34).
\]
Let
\[
 P=\gen{(13),(24)},\qquad Q=\gen{(14),(23)},\qquad
 D=\gen{a,b,r},
\]
and use the following six indexed subgroups:
\[
 A_1=\gen a,\quad A_2=\gen b,\quad A_3=P,\quad
 A_4=Q,\quad A_5=D,\quad A_6=\gen r.                 \tag{H1}
\]
Here $D$ has order eight and preserves the unordered partition
$\{\{1,2\},\{3,4\}\}$. Its normal forms are $a^i b^j r^\ell$, with
$i,j,\ell\in\{0,1\}$; conjugation by $r$ swaps $a$ and $b$.
The standard parabolics are exactly
\[
 \mathcal L=\{A,D,P,Q,\gen a,\gen b,\gen r,\gen s,1\}.       \tag{H2}
\]
Indeed, $D\cap P=\gen r$, $D\cap Q=\gen s$, and $P\cap Q=1$.
All four listed order-two groups lie in $D$; among them only $\gen r$
lies in $P$ and only $\gen s$ lies in $Q$. Distinct order-two groups
intersect trivially. These observations prove intersection closure and
show that every member of (H2) is a standard parabolic of (H1).

\begin{lemma}\label{N-lem:helly}
Every pairwise-intersecting family of left cosets of members of
$\mathcal L$ has a common representative.
\end{lemma}
\begin{proof}
Cosets of $A$ impose no restriction. A coset of $1$ is a singleton, so
pairwise intersection with it already gives the conclusion. Repeated
cosets can be removed; distinct cosets of the same subgroup are disjoint.
It remains to consider at most one coset of each of the other seven groups.

Call cosets of $\gen a,\gen b,\gen r,\gen s$ atomic. These are edges
with endpoints $x,xz$, where $z\in\{a,b,r,s\}$. Every atomic coset lies
in one left coset of $D$. Pairwise-intersecting atomic cosets therefore
lie in the same $D$-coset. The homomorphism
\[
 \epsilon:D\longrightarrow C_2,\qquad
 \epsilon(a^i b^j r^\ell)=i+j+\ell\pmod2
\]
takes all four of $a,b,r,s=abr$ to $1$. Thus their Cayley graph on each
$D$-coset is bipartite. Pairwise-intersecting edges in a triangle-free
simple graph have a common vertex: otherwise two edges meeting at $x$
would be met away from $x$ by a third edge, giving a triangle.

Suppose at least two atomic cosets occur, with common vertex $x$.
Their intersection is exactly $\{x\}$. A $D$-coset meeting an atomic
coset contains it, and hence contains $x$. Consider a $P$-coset meeting
all the atomic cosets. If the $\gen r$-coset occurs, it is contained in
that $P$-coset, so again $x$ belongs. Otherwise take two distinct atomic
generators from $a,b,s$. If the $P$-coset omitted $x$, it would contain
both opposite endpoints $xz_1,xz_2$, forcing $z_1z_2\in P$.
The product $ab=k$ is not in $P$, and $as,bs$ have order four, whereas
$P$ has exponent two. This is impossible. For $Q$ the same argument
uses $s$ in place of $r$ and the three remaining generators $a,b,r$.
Thus $x$ lies in every coset in the family.

Suppose exactly one atomic coset occurs. A $D$-coset is then redundant.
For the $\gen r$-coset a $P$-constraint is redundant, and for the
$\gen s$-coset a $Q$-constraint is redundant; the remaining pairwise
constraints have a common point. For an atom of type $a$ or $b$, the only
remaining possible failure would have a $P$-coset meeting one endpoint
and a $Q$-coset meeting the other. Their mutual intersection would imply
$z\in PQ$ (or $QP$) for $z=a$ or $b$. This does not happen. Multiplying
$a$ or $b$ by the four members $1,(14),(23),s$ of $Q$ gives, respectively,
the transposition $a$ or $b$, two three-cycles, and a four-cycle, none in
$P$. Inversion gives the same exclusion for $QP$.

Finally suppose no atomic coset occurs. Only a $P$-, a $Q$-, and a
$D$-coset remain. If all three occur, translate the unique intersection
of the first two to $1$, so those two cosets are $P,Q$. Write the third
as $gD$. It meets both only if $g^{-1}\in DP\cap DQ$. But
\[
 DP=D\cup D(13),\qquad DQ=D\cup D(14),\qquad DP\cap DQ=D.
\]
The two nontrivial right $D$-cosets are distinct, since $(13)(14)$ has
order three and cannot lie in the group $D$ of order eight. Hence
$g\in D$, so the third coset contains $1$ as well.
\end{proof}

Lemma~\ref{N-lem:helly} proves that every flag for (H1) has a common
representative $x$. It is therefore a translate by $x$ of a standard
flag and extends to $\{xA_i:1\leq i\leq6\}$. Thus the input is a
flag-transitive coset geometry, with no reliance on the finite enumeration.
Its Borel is $1$. Every minimal parabolic is $1$, since at least two of
the three distinct groups $\gen a,\gen b,\gen r$ remain when any single
type is deleted. Every standard panel consequently has one chamber,
so the geometry is not firm. It is not residually connected either:
the residue of the standard flag of types $\{1,5\}$ has group $\gen a$
and four remaining subgroups all equal to $1$. Its incidence graph is
two disjoint complete graphs on four vertices.

\subsection{Admissibility and the literal output}
Define $\phi(x)=bxb$. It fixes $A,D,\gen a,\gen b,1$ and interchanges
$P,Q$ and $\gen r,\gen s$. Hence it preserves exactly the lattice (H2)
and all intersections: it satisfies Definition 3.30. The named index
classes in Definition 3.31 are
\[
 \{1\},\quad\{2\},\quad\{3,4\},\quad\{5\},\quad\{6\}.    \tag{H3}
\]
In particular, $\phi(A_6)=\gen s$ is a standard parabolic but is not
one of the six named $A_i$.

For clarity, use the following coordinate convention for the infinite
group. Represent $xt^n$ by $(x,n)\in A\times\mathbb Z$ and put
\[
 (x,n)(y,m)=(x\phi^{-n}(y),n+m).                         \tag{H4}
\]
Associativity follows from the automorphism law, and
$(1,-1)(x,0)(1,1)=(\phi(x),0)$. Conversely the relation in the HNN
presentation moves every letter from $A$ to the left, giving $xt^n$;
the map into (H4) shows uniqueness. Thus (H4) is precisely the presented
HNN extension, not merely a quotient of it. In this example
$\phi^{-n}=\phi^n$ because $\phi^2=1$.

Let $V=\gen{r,s}=\{1,r,s,k\}$. For each $\phi$-invariant subgroup
$B\leq A$, the subgroup $\gen{B,t}$ is the set $B\times\mathbb Z$
under (H4). The literal maximal parabolics of the output are therefore
\[
\begin{array}{c|c}
\text{type}&\text{subgroup in coordinates (H4)}\\\hline
\{1\}&\gen a\times\mathbb Z\\
\{2\}&\gen b\times\mathbb Z\\
\{3,4\}&\{1\}\times\mathbb Z\\
\{5\}&D\times\mathbb Z\\
\{6\}&V\times\mathbb Z\\
t&A\times\{0\}
\end{array}                                                     \tag{H5}
\]
Here the sixth input group expands because conjugation by $t$ changes
$r$ to $s$, and $P\cap Q=1$ accounts for the third row.

\begin{theorem}
The output of Definition 3.31 for (H1) and $\phi(x)=bxb$ is neither
flag-transitive nor a geometry.
\end{theorem}
\begin{proof}
Write $G_1,G_2,G_6$ for the output groups at the corresponding singleton
classes. The three distinct-type cosets
\[
 G_1,\qquad G_2,\qquad aG_6                              \tag{H6}
\]
meet pairwise: their respective intersections contain $1,a,b$, since
$ab=k\in V$. On the other hand,
\[
 G_1\cap G_2=\{1\}\times\mathbb Z=\gen t,
 \qquad aG_6=aV\times\mathbb Z.
\]
The coset $aV$ omits $1$ (its elements are odd permutations), so (H6)
has empty total intersection. A translate of the standard flag has a
common representative. Hence this flag cannot be in its orbit, proving
failure of flag transitivity.

More strongly, the output type $\{3,4\}$ has subgroup $\gen t$.
Any one of its cosets incident with both $G_1$ and $G_2$ is contained
in both: $\gen t$ is a subgroup of each. Such a coset must be exactly
$\gen t$, which does not meet $aG_6$. Thus (H6) cannot extend to a
chamber. Since the type set is finite, it extends to a maximal flag
that is not a chamber. The output is not a geometry.
\end{proof}


\subsection{The simplex HNN construction}
Take $a_i=(i,i+1)$ in $S_6$, $1\leq i\leq5$, as in
Example 3.36. The associated subgroups are
$C=\langle a_2,a_3\rangle$ and $D=\langle a_3,a_4\rangle$,
with indexed map $a_2\mapsto a_3$, $a_3\mapsto a_4$.
The type classes are $\{1\},\{2,3,4\},\{5\}$ and $F=\{4\}$.
The indexed HNN theorem in the common paper therefore gives
\[
 G_1=\langle A_1,t\rangle,\quad G_{234}=\langle a_1,a_5,t\rangle,
 \quad G_5=\langle A_5,t\rangle,\quad
 G_t=A_4=\langle a_1,a_2,a_3,a_5\rangle.
\]
Indexed compatibility follows from the named bijection and classical
Coxeter parabolic embeddings. The resulting geometry is
flag-transitive, residually connected and firm. Its Borel is trivial,
and the exact panel formula gives $2,6,2,\infty$.
Example 3.37's whole-base automorphism gives $A\rtimes\ZZ$ with
the stated orbit-coarsened parabolics when its indexed action is
specified. For firm inputs the Boolean-parabolic lemma in the common paper
supplies the indexed action from source admissibility. Preservation
of an unlabelled parabolic family alone does not supply it on an
arbitrary nonfirm input.
